%--------------------------------------------------------------------------------
%
%
%--------------------------------------------------------------------------------
\pagebreak
\section*{ADD — Integer Addition}
\addcontentsline{toc}{section}{ADD - Integer Addition}

%--------------------------------------------------------------------------------
\instructionentry{Syntax}{

  \texttt{ADD RegR, RegB, RegA} \\
  \texttt{ADD RegR, RegB, val}
  
  \texttt{ADD[.D/W/H/B] RegR, ofs ( RegB )} \\
  \texttt{ADD[.D/W/H/B] RegR, RegX ( RegB )}
}

%--------------------------------------------------------------------------------
\instructionentry{Format}{

The ADD instruction uses the register, the immediate-15, the indexed and the register indexed instruction formats. \\
	
\begin{flushleft}	

	\drawInstrFormat
  	{0,1,3,5,6.5,8.5,9.5,11.5,16}
  	{\OpCodeGrpALU,\OpCodeADD,RegR, 0, RegB, 0, RegA, 0}
	
	\drawInstrFormat
  	{0,1,3,5,6.5,8.5,16}
  	{\OpCodeGrpALU,\OpCodeADD, RegR, 1, RegB, val}

	\drawInstrFormat
  	{0,1,3,5,6.5,8.5,9.5,16}
  	{\OpCodeGrpMEM,\OpCodeADD, RegR, 0, RegB, dw, ofs}
 
	\drawInstrFormat
  	{0,1,3,5,6.5,8.5,9.5,11.5,16}
  	{\OpCodeGrpMEM,\OpCodeADD, RegR, 1, RegB, dw, RegX, 0}

\end{flushleft}  
}

%--------------------------------------------------------------------------------
\instructionentry{Description}{

The \texttt{ADD} instruction adds two signed 64-bit operands and stores the result to the target register \texttt{RegR}. The instruction supports the immediate, the register and the two memory reference instruction formats. An arithmetic overflow results in a numeric overflow trap. The indexed instruction formats may cause a memory reference associated trap.
}

%--------------------------------------------------------------------------------
\instructionentry{Operation}{

  	\texttt{RegR <- RegB + RegA} {(register format)}\\
  	\smallskip
  	\texttt{RegR <- RegB + immOperand( )} {(immediate format)}\\
   	\smallskip
  	\texttt{RegR <- RegR + memOperand( )} {(indexed formats)}\\
}

%--------------------------------------------------------------------------------
\instructionentry{Exceptions}{

 	Integer Overflow Trap\\
	Data Alignment Trap\\
	Memory Access Violation Trap\\
	Memory Protection Violation Trap\\
	Data TLB Miss Trap
}

%--------------------------------------------------------------------------------
\instructionentry{Notes}{

	None. 
}
